\documentclass[10pt]{article}
\usepackage[paperwidth=170mm,paperheight=240mm,top=22mm,bottom=22mm,
  inner=20mm,outer=20mm,headsep=6mm]{geometry}
\usepackage[T1]{fontenc}
\usepackage{lmodern}
\usepackage{amsmath,amssymb}
\usepackage{microtype}
\usepackage{fancyhdr}
\usepackage[hidelinks]{hyperref}
\usepackage{xcolor}

\hyphenation{multi-par-ti-tion}

\usepackage{tikz}

% ------------------------------------------------------------ FM-like layout
\newif\iffmlabels \fmlabelsfalse  % \fmlabelstrue shows the Mizar labels
\newcounter{fmsec}
% keep a heading with the text that follows it
\newcommand{\fmkeep}{\par\vskip 0pt plus 8\baselineskip\penalty-200
  \vskip 0pt plus -8\baselineskip}
\newcommand{\fmsection}[1]{\stepcounter{fmsec}%
  \fmkeep\bigskip\begin{center}\normalsize\scshape\arabic{fmsec}.\ #1\end{center}\nopagebreak
  \smallskip}
\newcommand{\fmunnumbered}[1]{%
  \fmkeep\bigskip\begin{center}\normalsize\scshape #1\end{center}\nopagebreak\smallskip}
\newcommand{\fmpara}[1]{\par\smallskip #1\par}
\newlength{\fmnum}\setlength{\fmnum}{2.6em}
\newlength{\fmdefnum}\setlength{\fmdefnum}{4.6em}
\makeatletter
% (n) statement (and, optionally, the Mizar label at the right)
\newcommand{\fmthm}[3]{\par\smallskip\noindent
  \hangindent=\dimexpr\fmnum+1.2em\relax\hangafter=1%
  \hspace*{1.2em}\makebox[\fmnum][l]{(#1)}#2%
  \iffmlabels\nobreak\hfill\nobreak
    \hbox{\hspace{1em}\scriptsize\textsf{\color{black!60}#3}}\fi
  \phantomsection\def\@currentlabel{#1}\label{thm:#3}\par}
\makeatother
\newcommand{\fmdef}[2]{\par\smallskip\noindent
  \hangindent=\fmdefnum\hangafter=1%
  \makebox[\fmdefnum][l]{(Def.\,#1)}#2\par}
\newcounter{fmi}
\newif\iffmsameline
% \fmsameline: the first item goes on the line of the theorem number
\newcommand{\fmsameline}{\global\fmsamelinetrue}
\newenvironment{fmlist}{\setcounter{fmi}{0}%
  \renewcommand{\item}{\stepcounter{fmi}%
    \iffmsameline\global\fmsamelinefalse\else
      \par\noindent\hspace*{\dimexpr\fmnum+1.2em\relax}\fi
    \hangindent=\dimexpr\fmnum+1.2em+2.8em\relax\hangafter=1%
    \makebox[2.8em][r]{(\roman{fmi})\ }}%
  \iffmsameline\else\par\fi}{\par}
% the period that ends a statement goes into its last item
\def\fmlistend#1{#1\end{fmlist}}
% text after a list, aligned with the statement
\newcommand{\fmafterlist}[1]{\par\noindent
  \hangindent=\dimexpr\fmnum+1.2em\relax\hangafter=0 #1}
\newcommand{\thref}[1]{(\ref{thm:#1})}
% the ORCID iD icon
\newcommand{\orcidicon}{\tikz[baseline=-0.6ex]{
  \fill[green!55!black!80!lime] (0,0) circle (0.9ex);
  \node[white,font=\fontsize{4.5}{4.5}\selectfont\sffamily\bfseries,inner sep=0]
    at (0,0) {iD};}}
\newcommand{\fmbox}[1]{\fcolorbox{cyan!70}{white}{\strut #1}}
\setlength{\fboxsep}{1.5pt}\setlength{\fboxrule}{0.4pt}
\setlength{\parindent}{1.5em}

\pagestyle{fancy}
\fancyhf{}
\fancyhead[LE]{\small\thepage}
\fancyhead[RO]{\small\thepage}
\fancyhead[CE]{\small\scshape daniel severín}
\fancyhead[CO]{\small\scshape Additive chromatic number of several families of graphs}
\renewcommand{\headrulewidth}{0pt}
\fancypagestyle{first}{\fancyhf{}\renewcommand{\headrulewidth}{0pt}
  \fancyfoot[C]{\small\thepage}}

\begin{document}
\thispagestyle{first}

\noindent{\scriptsize\scshape A Mizar article in the style of
Formalized Mathematics}
\vspace{16mm}

\begin{center}
{\LARGE\bfseries On the Additive Chromatic Number of\\[2pt] Several Families of Graphs}
\par\vspace{12mm}
{\large Daniel Severín\,\href{https://orcid.org/0000-0003-2391-3489}{\orcidicon}\\
Facultad de Ciencias Exactas, Ingeniería y Agrimensura\\
Universidad Nacional de Rosario\\
Argentina}\\[3pt]
{\small\href{mailto:daniel@fceia.unr.edu.ar}{\texttt{daniel@fceia.unr.edu.ar}}}
\end{center}

\vspace{10mm}
\begin{list}{}{\leftmargin=8mm\rightmargin=8mm}\item\relax
{\small
\textbf{\normalsize Summary.} Given a graph $G$, a labeling $f$ of its vertices
with natural numbers is an additive coloring of $G$ if the sums $f(N(u))$ and
$f(N(v))$ of the labels in the open neighborhoods of $u$ and $v$ differ for
every edge $uv$ of $G$. The additive chromatic number $\eta(G)$ is the least
$k$ such that $G$ has an additive coloring with labels in $\{1,\ldots,k\}$
\cite{czerwinski}. We formalize Sections 1--5 of \cite{severin}, including the
proofs of its propositions that are given in the appendix of its arXiv
version: $\eta(G)$ of regular bipartite, complete multipartite and complete
split graphs, of joins with complete graphs, fans, windmills, wheels, odd
cycles, headless spiders and three families of suns, and the Additive
Coloring Conjecture $\eta(G) \leq \chi(G)$ for complete multipartite graphs,
for joins with complete graphs and for split graphs. The two results that
the paper cites from \cite{marenco} (the value of the topological additive
number of complete multipartite digraphs and a lower bound from cliques), and
the equality $\eta(G) = \min \eta_t(D)$ over the acyclic orientations $D$ of
$G$, are proved as well.

\par\medskip
MSC: \href{https://zbmath.org/classification/?q=05C15}{\fmbox{05C15}}\enspace
\href{https://zbmath.org/classification/?q=05C78}{\fmbox{05C78}}\enspace
\href{https://zbmath.org/classification/?q=68V20}{\fmbox{68V20}}

\smallskip
Keywords: additive coloring;\enspace lucky labeling;\enspace additive
chromatic number;\enspace split graph;\enspace complete multipartite graph

\smallskip
MML identifier: \fmbox{\texttt{ACN}},\enspace version:
\fmbox{\texttt{8.1.15 5.94.1493}}
}
\end{list}
\vspace{4mm}

\fmunnumbered{Introduction}

The Additive Coloring Problem asks for a labeling $f : V(G) \to
\{1,\ldots,k\}$ such that $f(N(u)) \neq f(N(v))$ for every edge $uv$, with
$k$ as small as possible; the least such $k$ is the additive chromatic number
(or lucky number) $\eta(G)$ \cite{czerwinski}. The Additive Coloring
Conjecture states that $\eta(G) \leq \chi(G)$ for every graph $G$. This
article formalizes the paper \cite{severin} in the Mizar system
\cite{mizar,fourdecades}. The proofs follow an earlier formalization of the
same paper in Isabelle/HOL.

Graphs are finite simple graphs given as non empty, finite, irreflexive and
symmetric relational structures \cite{necklace}; $\Omega_G$ is the set of
vertices of $G$ and $E_G$ its (symmetric) adjacency relation. The chromatic
number $\chi(G)$ is the least number of classes of a partition of $\Omega_G$
into stable sets, as defined in \cite{mycielsk,dilworth}. Sums of labels
over finite sets are sums of bags \cite{uproots,nat5}. The additive chromatic
number is introduced in Section~2, together with acyclic orientations
(based on the reduction sequences of \cite{rewrite1}), topological additive
numberings, the topological additive number $\eta_t(D)$ of \cite{marenco}
(an extended real, $+\infty$ when there is no numbering) and the equality
$\eta(G) = \min \eta_t(D)$ \thref{EtaMinEtaT}.

Each family of graphs of \cite{severin} is introduced as a predicate on
graphs, with a theorem stating that such graphs exist. The results of the
paper correspond to the following theorems: Observation~1
\thref{Observation1}, Lemma~1 \thref{Lemma1}, Corollary~1
\thref{Corollary1}, Lemma~2 \thref{Lemma2Value}, \thref{Lemma2Finite},
Theorem~1 \thref{Theorem1}, \thref{Theorem1Conjecture}, Observation~2
\thref{Observation2}, Theorem~2 \thref{Theorem2},
\thref{Theorem2Conjecture}, fans \thref{Fan}, windmills \thref{Windmill},
Proposition~1 \thref{Proposition1}, wheels \thref{Wheel}, Theorem~3
\thref{Theorem3}, \thref{Theorem3Conjecture}, \thref{SplitConjecture},
Proposition~2 \thref{Proposition2}, Lemma~3 \thref{Lemma3}, Corollary~2
\thref{Corollary2}, Proposition~3 \thref{Proposition3Thin},
\thref{Proposition3Thick}, Proposition~4 \thref{Proposition4} and
Proposition~5 \thref{Proposition5Cycle}, \thref{Proposition5Wheel}. The
labeling of complete suns is the one of the appendix of \cite{severin}, with
one correction: in the case $m \equiv 2 \pmod 6$, $i = \lfloor m/2\rfloor +
1$, the label of $v_{i-1}$ is $2$.

Throughout, $N(v)$ is the open neighborhood of $v$, $\deg v$ its degree,
$f(N(v))$ the sum of the labels in $N(v)$, $\overline{\overline{X}}$ the
cardinality of $X$, $\lceil a/b\rceil$ the ceiling of $a/b$, and
$\mathrm{succ}_m(i)$, $\mathrm{pred}_m(i)$ the successor and the predecessor
of $i$ modulo $m$ in $\{1,\ldots,m\}$. In the families on two sides
(spiders and suns), $u_i = \langle 0, i\rangle$ and $v_i = \langle 1,
i\rangle$.

\input{body}

\bigskip
\noindent\textsc{Acknowledgement:} This document has the layout of a paper in
\emph{Formalized Mathematics}. It is generated by \texttt{miz2tex.py} from the
Mizar article \texttt{acn.miz} and has not been published.

\makeatletter
{\footnotesize
\fmunnumbered{References}
\let\section\@gobbletwo
\setlength{\itemsep}{0pt}
\input{refs}
}
\makeatother

\end{document}
